\section{The Recursive Spawning Protocol}
\label{sec:rs-protocol}

The Recursive Spawning Protocol (RSP) is the structural mechanism that converts a delegation chain from a forensic reconstruction assembled after the fact from mutable logs and policy assertions into a runtime-verified artifact enforced at every spawn event without exception. RSP operates across all three \bxthree{} layers in mandatory concert: the Purpose Layer authorizes, the Bounds Engine constrains, and the Fact Layer records and enforces. There are no advisory phases, no configurable bypasses, and no operational circumstances under which any machine actor may circumvent any protocol requirement.

The protocol is defined for a system $\mathcal{S} = \langle \mathcal{N}, \alpha, R, h, \mathcal{L}, D_{\max} \rangle$ where $\mathcal{N}$ is the agent node set, $\alpha(n)$ is the immutable parent pointer, $R(n)$ is the authorized operating scope, $h(n)$ is the human anchor, $\mathcal{L}$ is the Fact Layer forensic ledger, and $D_{\max}$ is the architectural maximum spawn depth. Every spawn event executes through five mandatory phases in strict sequence; failure at any phase terminates the spawn attempt.

% ============================================================================
\subsection{Phase 1: Purpose Layer Validation}
\label{sec:rs-phase1}

Every spawn event begins with a Purpose Layer validation gate. The Purpose Layer is the exclusive authorization origin for all chain growth — no machine actor may initiate a spawn without a matching Purpose Layer warrant recorded in the Fact Layer forensic ledger before execution. This is enforced structurally by the Fact Layer pre-commit hook, which rejects any provisioning request lacking a matching warrant before the operation reaches the ledger write stage.

\medskip
\noindent\textbf{Validation conditions.} A parent node $n_i \in \mathcal{N}$ intending to spawn a child $n_{i+1}$ must submit a spawn request declaring: (P1) the proposed child operating scope $R(n_{i+1})$, (P2) the specific operational condition necessitating spawning, and (P3) a demonstration that the condition cannot be resolved within $R(n_i)$. The Purpose Layer validates three mandatory conditions:

\begin{enumerate}
    \item[(P1)] \textbf{Scope refinement.} $R(n_{i+1})$ must be a strict descendant refinement of $R(n_i)$:
    \begin{equation}
    R(n_{i+1}) \subset R(n_i).
    \label{eq:rs-scope-refinement}
    \end{equation}
    Lateral scope stability — where $R(n_{i+1}) = R(n_i)$ — is not authorized. The chain grows through scope \emph{narrowing}, not sideways expansion.

    \item[(P2)] \textbf{Operational necessity.} The parent must declare the precise condition requiring child spawning and identify why it cannot be resolved within $R(n_i)$.

    \item[(P3)] \textbf{Minimum scope proportionality.} The child scope must be the narrowest scope sufficient to address the declared condition. Scope beyond minimum necessary blocks warrant issuance.
\end{enumerate}

\medskip
\noindent\textbf{Warrant issuance.} When all three conditions are satisfied, the Purpose Layer issues a cryptographically signed spawn warrant $w_i$ written atomically to the Fact Layer ledger:

\begin{equation}
w_i = \bigl\langle
\text{parent}=n_i,\;
\text{child}=n_{i+1},\;
R(n_{i+1}),\;
\text{justification},\;
\text{timestamp},\;
\sigma_{\text{PL}}
\bigr\rangle .
\label{eq:rs-warrant}
\end{equation}

The warrant $\sigma_{\text{PL}}$ is the sole authorized precondition for spawn. No machine actor may initiate a spawn without a matching ledger entry. At root provisioning, the Purpose Layer records $h(n_0) = h$ in the ledger. The anchor is non-collapsing: for all nodes $n_j$ descended from $n_0$, $h(n_j) = h(n_0)$. No machine actor can serve as a chain root.

% ============================================================================
\subsection{Phase 2: Bounds Engine Depth Check}
\label{sec:rs-phase2}

When a Purpose Layer warrant is confirmed, the spawn request proceeds to the Bounds Engine for depth evaluation. The Bounds Engine enforces the depth bound as a structural invariant — not a policy threshold overrideable by operational urgency.

The system enforces a fixed maximum chain depth $D_{\max}$, defined by the Purpose Layer at system initialization and recorded in the Fact Layer authorization ledger. $D_{\max}$ is not reconfigurable at runtime by any machine actor. The depth of any node $n$ is:

\begin{equation}
\text{depth}(n) =
\begin{cases}
0 & \text{if } \alpha(n) = \bot \quad (\text{root human anchor}) \\[6pt]
1 + \text{depth}(\alpha(n)) & \text{otherwise}.
\end{cases}
\label{eq:rs-depth}
\end{equation}

At each spawn event $n_i \xrightarrow{\text{spawn}} n_{i+1}$, the Bounds Engine evaluates:
\begin{equation}
\text{depth}(n_i) + 1 \leq D_{\max}.
\label{eq:rs-depth-check}
\end{equation}

If Equation~\ref{eq:rs-depth-check} fails, the Bounds Engine rejects the spawn, logs the rejection as a ledger entry, and transitions the parent to \texttt{ESCALATING} state. The refusal is deterministic: there is no override, no emergency pathway, and no urgency flag that permits exceeding $D_{\max}$.

The Purpose Layer additionally defines $D_{\text{ESCALATE}}$, where $0 < D_{\text{ESCALATE}} \leq D_{\max}$, as the depth at which human confirmation is required before chain growth continues — even when the spawn is technically within $D_{\max}$. When $\text{depth}(n_i) \geq D_{\text{ESCALATE}}$, the Bounds Engine suspends autonomous execution until the human anchor $h(n_i)$ issues \texttt{ACK}, \texttt{RESOLVE}, or \texttt{REJECT}. No machine actor may issue any of these directives, and no machine actor may serve as the terminus of an RSP chain.

% ============================================================================
\subsection{Phase 3: Fact Layer — Immutable Pointer and Forensic Ledger}
\label{sec:rs-phase3}

When the Purpose Layer warrant is confirmed and the Bounds Engine depth check passes, the spawn request reaches the Fact Layer, which provides the structural enforcement that makes all RSP guarantees runtime-operative.

The Fact Layer creates child node $n_{i+1}$ with an immutable parent pointer $\alpha(n_{i+1}) = n_i$. This pointer satisfies:

\begin{equation}
\forall n,\; p,\; c:\;
\alpha(c) = p \;\land\; \text{provisioned}(c)
\;\implies\;
\neg\,\exists\,\text{modify}(\alpha(c)).
\label{eq:rs-pointer-immutability}
\end{equation}

No actor — including the Purpose Layer after provisioning, the Bounds Engine, any administrative process, or any external principal — may overwrite, redirect, or delete $\alpha(n_{i+1})$ after the provisioning write is confirmed. The immutability is protocol-enforced, not access-control-enforced. In access-control models, immutability can be circumvented by sufficient privilege escalation. In the Fact Layer write protocol, the modification operation is structurally impossible regardless of credential compromise or software vulnerability. There is no administrative override and no privileged actor capable of altering the chain after provisioning.

The Fact Layer records every authorized spawn event as an append-only ledger entry:

\begin{equation}
\ell_j = \bigl\langle
\text{index}=j,\;
\text{node}=n_i,\;
\text{event-type},\;
\text{payload},\;
\text{timestamp},\;
\text{attribution},\;
\text{hash}(\ell_{j-1})
\bigr\rangle .
\label{eq:rs-ledger-entry}
\end{equation}

The hash-chaining establishes continuity: no entry can be inserted, deleted, reordered, or altered without breaking chain continuity detectable by any observer with ledger read access. The complete spawn topology — warrant, parent pointer, and child configuration — is committed atomically. There is no temporal window in which the parent pointer exists without a matching Purpose Layer warrant.

The Fact Layer pre-commit hook verifies two structural invariants before any spawn is recorded:

\begin{enumerate}
    \item[(FC1)] $R(n_{i+1}) \subseteq R(n_i)$ — enforced as a structural invariant.
    \item[(FC2)] $\text{depth}(n_{i+1}) \leq D_{\max}$ — re-verified at the Fact Layer even though checked at the Bounds Engine.
\end{enumerate}

The conjunction of (FC1) and (FC2) at every spawn step makes the Scope Refinement Invariant a structural guarantee for the lifetime of the chain.

% ============================================================================
\subsection{Phase 4: Child Activation with Immutable Parent}
\label{sec:rs-phase4}

When the Fact Layer pre-commit hook confirms both constraints and records the spawn event, the child node $n_{i+1}$ activates with three structurally immutable properties established at provisioning:

\begin{itemize}
    \item \textbf{Immutable parent reference.} $\alpha(n_{i+1}) = n_i$ is a Fact Layer property of the child. The child has read access to its parent reference but no mechanism to modify it. A terminated, crashed, or decommissioned parent does not affect the child's ability to reference its parent through $\alpha(\cdot)$.
    \item \textbf{Authorized operating scope.} $R(n_{i+1})$ is established by the Purpose Layer spawn warrant and recorded in the Fact Layer ledger. The child operates strictly within $R(n_{i+1})$ — any action outside triggers the exception handling protocol. The child cannot expand $R(n_{i+1})$ without a new Purpose Layer warrant.
    \item \textbf{Human anchor inheritance.} The child inherits $h(n_{i+1}) = h(n_0)$ from its parent, established at root provisioning. This reference is identical for all nodes in the chain and immutable at every node.
\end{itemize}

% ============================================================================
\subsection{Phase 5: Chain Verification}
\label{sec:rs-phase5}

RSP provides a structural chain verification procedure executable at any runtime point. To verify that a node $n$ is authorized by a human principal, an auditor performs:

\begin{enumerate}
    \item[(V1)] \textbf{Traverse the immutable parent pointer chain.} Repeatedly read $\alpha(n)$, then $\alpha(\alpha(n))$, until reaching $\alpha(n) = \bot$. Record $C(n) = \langle n, \alpha(n), \alpha(\alpha(n)), \ldots, n_0 \rangle$.
    \item[(V2)] \textbf{Verify the human anchor.} Confirm $n_0$ has $h(n_0)$ recorded and $\alpha(n_0) = \bot$.
    \item[(V3)] \textbf{Verify Purpose Layer warrants.} For each link $\alpha(n_{j+1}) = n_j$, confirm the spawn event was recorded in the forensic ledger with a matching warrant from Equation~\ref{eq:rs-warrant}.
    \item[(V4)] \textbf{Verify scope refinement at each link.} Confirm $R(n_{j+1}) \subseteq R(n_j)$ for each consecutive pair.
    \item[(V5)] \textbf{Verify depth bound.} Confirm $\text{depth}(n_j) \leq D_{\max}$ for every node in the chain.
\end{enumerate}

If all five steps pass, the chain is verified. The verification outcome is identical regardless of which node initiates it — there is no self-serving interpretation of authorization that a drifted or orphaned node could produce. The chain terminates at the human anchor $h$ by architectural constraint. The human anchor is non-collapsing and non-substitutable. No machine actor can serve as the terminus of an RSP chain.

% ============================================================================
\subsection{Depth Management: Max Depth, Spawn Refusal, and Convergence}
\label{sec:rs-depth-mgmt}

The combination of $D_{\max}$, spawn refusal at the depth limit, and convergence criteria constitutes the complete depth management system of RSP — three components of a single depth enforcement architecture ensuring every RSP chain terminates within a finite number of steps, at or before the architectural depth bound, and always at a human principal.

\medskip
\noindent\textbf{Maximum spawn depth $D_{\max}$.} Set by the Purpose Layer at system initialization based on the operational consequence domain. High-stakes environments require lower $D_{\max}$ values. The value is recorded in the Fact Layer ledger and is immutable. The depth bound is enforced at every spawn step by both the Bounds Engine and the Fact Layer pre-commit hook.

\medskip
\noindent\textbf{Spawn refusal at the depth limit.} When $\text{depth}(n_i) + 1 > D_{\max}$, the Bounds Engine refuses the spawn and enters \texttt{ESCALATING} state. The refusal is deterministic: the Bounds Engine has no discretionary authority to exceed $D_{\max}$ regardless of operational urgency or task criticality. The refusal is logged and the human anchor is notified. The chain does not proceed, does not bypass the depth bound, and does not continue autonomously past $D_{\max}$. The $D_{\text{ESCALATE}}$ threshold provides advance warning — human confirmation is required before the depth limit is reached, not only after it is breached.

\medskip
\noindent\textbf{Convergence criteria.} A chain has \emph{converged} when all of the following hold simultaneously:

\begin{enumerate}
    \item[(C1)] \textbf{Terminal state reached.} All leaf nodes are in \texttt{COMPLETED}, \texttt{TERMINATED}, or \texttt{ESCALATING} state with a \texttt{REJECT} directive from $h$.
    \item[(C2)] \textbf{No active exceptions.} No exception condition remains unresolved within the chain's operating scope.
    \item[(C3)] \textbf{Ledger continuity maintained.} All state transitions and spawn events are recorded in the Fact Layer forensic ledger with hash-chain continuity intact.
    \item[(C4)] \textbf{Verified human termination.} The chain terminates at the human anchor $h$ — the root has a human principal designation and $\alpha(n_0) = \bot$.
\end{enumerate}

Convergence is assessed at each state transition event. When all four criteria hold, the chain is \emph{resolved}: no further autonomous spawn events are authorized, and the chain is archived as a completed ledger entry. Convergence is not merely the absence of active nodes — it is a positive confirmation that the chain has reached a defined terminal state with a complete, tamper-evident accountability record tracing to a human principal.

\medskip

The Recursive Spawning Protocol provides a single architectural guarantee: that every agent in a multi-agent system operates within a delegation chain that is simultaneously immutable, auditable, scope-refined, depth-bounded, and verified to terminate at a human principal. The guarantee is structural — it holds regardless of the goodness, reliability, or intentions of any machine actor in the chain. This is the operational expression of the upstream \bxthree{} accountability guarantee applied to recursive agent population dynamics.